\documentclass[11pt,a4paper]{article}
\usepackage{apnotes}
\setsunday{3}
\begin{document}

\notestitle{Transformers \& LLM Engineering}{11--15}{Understand a GPT from the inside, build one from scratch, teach it a new instruction with LoRA, evaluate it honestly and serve it with structured output --- project \textbf{P3 --- MiniLLM}.}

\section*{How to use these notes}
These notes follow the Sunday 3 slides topic by topic and explain each step in full, using the same colour code:
\legendbox
Numbers in \emph{Our notebook result} boxes come from \texttt{notebooks/S03\_Transformers\_LLM\_Engineering.ipynb}. We build every model ourselves so it runs on a CPU; the same ideas scale to open models with Hugging Face \texttt{transformers} and \texttt{peft} on a GPU. The last section is your \textbf{Monday--Saturday project homework}.

\begin{tabularx}{\linewidth}{@{}c c X l@{}}
\toprule
\textbf{Hour} & \textbf{Mode} & \textbf{Focus} & \textbf{Dataset}\\
\midrule
11 & \teachtag & Language models, tokens, attention, GPT from scratch & \ds{Tiny Shakespeare}\\
12 & \teachtag & Encoders, transfer learning, strong baselines & \ds{SST-2}\\
13 & \buildtag & Instruction tuning with LoRA & \ds{SST-2 instructions}\\
14 & \buildtag & Evaluation, contamination, forgetting & \ds{SST-2 test}\\
15 & \buildtag & Decoding, structured output, serving & \ds{Prompt suite}\\
\bottomrule
\end{tabularx}

\begin{bdef}[Learning objectives]
By the end of today you can (1) define a language model and perplexity; (2) explain tokenisation (BPE) and embeddings; (3) derive and compute scaled dot-product attention; (4) build a GPT and fine-tune it with LoRA; (5) evaluate an LLM honestly --- baselines, contamination, forgetting; (6) serve it with sampling, structured output and batching.\\[2pt]
\textbf{Prerequisites:} Sunday 2 --- neural networks, softmax, cross-entropy, backpropagation, PyTorch.
\end{bdef}

%% =====================================================================
\hourhead{11}{TEACH}{Language models, tokens and attention}{Tiny Shakespeare (1.1\,M characters)}{Attention maths, tokenisers, GPT internals}

\topic{What is a language model?}
\begin{bdef}[Language model]
A model that gives the probability of the next token given all previous tokens. Generating text means repeatedly sampling a next token and appending it.
\end{bdef}
\begin{bmath}[Chain rule of probability, loss and perplexity]
$P(x_1,\dots,x_T)=\prod_{t=1}^{T}P(x_t\mid x_{<t})$\\[3pt]
Training loss (average negative log-likelihood): $\mathcal{L}=-\frac1T\sum_t\log P(x_t\mid x_{<t})$\\[3pt]
Perplexity $=e^{\mathcal{L}}$ --- roughly, how many choices the model is torn between at each step.
\end{bmath}
The chain rule shows why next-token prediction is enough: if a model knows the probability of every next token, it knows the probability of any whole text. Training simply maximises the probability of real text, which is the same as minimising the cross-entropy loss from Sunday 2, applied at every position.
\begin{bex}[A baseline to beat]
With 65 possible characters, uniform guessing has perplexity 65. After the line ``ROMEO:'' and a line break, a good model concentrates probability on capital letters, so its perplexity is far lower.
\end{bex}

\topic{Tokenisation and embeddings}
\begin{bdef}[Token]
The unit a model reads: a character, a whole word, or a sub-word piece. Modern LLMs use sub-word tokens.
\end{bdef}
Characters give a tiny vocabulary but very long sequences; whole words give short sequences but cannot handle new words. Sub-word tokenisation is the compromise: common words are single tokens, rare words are split into pieces, and nothing is ever ``unknown''.
\begin{balgo}[Byte-pair encoding (BPE)]
\begin{algorithmic}[1]
\State start with single characters plus an end-of-word mark
\For{$k=1$ to $K$ merges}
  \State count every adjacent pair of symbols in the corpus
  \State merge the most frequent pair into a new symbol
\EndFor
\State \Return the list of merges (this defines the vocabulary)
\end{algorithmic}
\end{balgo}
\begin{bres}
The first merges learned on Shakespeare are \texttt{e</w>}, \texttt{s</w>}, \texttt{th}, \texttt{t</w>}, \texttt{d</w>}, \ldots, \texttt{the</w>} --- word endings and the most common words appear first.
\end{bres}
\begin{bdef}[Embedding]
A learned vector for each token: $e_i=E[i]$ with $E\in\mathbb{R}^{V\times d}$. Training moves tokens used in similar contexts to nearby vectors.
\end{bdef}

\topic{Scaled dot-product attention}
\begin{bdef}[Attention]
Each token produces a \emph{query} (what it is looking for), a \emph{key} (what it offers) and a \emph{value} (the information it passes on). A token gathers the values of the tokens whose keys best match its query.
\end{bdef}
\begin{bmath}[Formula]
$Q=XW_Q,\;K=XW_K,\;V=XW_V$ \qquad $\mathrm{Attention}(Q,K,V)=\mathrm{softmax}\!\Big(\dfrac{QK^\top}{\sqrt{d_k}}+M\Big)V$\\[3pt]
Causal mask: $M_{ij}=0$ if $j\le i$ and $-\infty$ if $j>i$, so a token cannot see the future.
\end{bmath}
Dividing by $\sqrt{d_k}$ keeps the dot products from growing with the vector size; without it, the softmax saturates (one weight near 1, the rest near 0) and gradients vanish.
\begin{bex}[Two tokens by hand]
Token 2 has scores $q_2\cdot k_1=1$ and $q_2\cdot k_2=3$, with $d_k=4$. Scaled scores: $0.5$ and $1.5$. Softmax: $0.27$ and $0.73$. Token 2's output is $0.27\,v_1+0.73\,v_2$.
\end{bex}
\begin{bres}
In the notebook's 5-token example every row of attention weights sums to 1 and every entry above the diagonal is exactly 0.
\end{bres}

\topic{The Transformer block and positions}
\begin{balgo}[GPT forward pass]
\begin{algorithmic}[1]
\State $h\gets$ token embedding $+$ position embedding
\For{each of $L$ blocks}
  \State $h\gets h+\mathrm{MultiHeadAttention}(\mathrm{LayerNorm}(h))$ \Comment{tokens exchange information}
  \State $h\gets h+\mathrm{MLP}(\mathrm{LayerNorm}(h))$ \Comment{each token processed on its own}
\EndFor
\State logits $\gets\mathrm{LayerNorm}(h)\,W_{out}$ \Comment{one score per vocabulary entry}
\end{algorithmic}
\end{balgo}
\begin{bdef}[Multi-head attention, residuals, LayerNorm]
\textbf{Multi-head}: several attentions in parallel, each free to track a different pattern (speaker names, rhyme, grammar). \textbf{Residual connections} ($h+\dots$) let gradients flow through deep stacks. \textbf{LayerNorm} keeps activations at a stable scale.
\end{bdef}
\begin{bmath}[Sinusoidal positions (original Transformer)]
$PE_{(pos,2i)}=\sin\!\big(pos/10000^{2i/d}\big)$, \;$PE_{(pos,2i+1)}=\cos\!\big(pos/10000^{2i/d}\big)$. Our MiniGPT learns its position embeddings instead.
\end{bmath}
\begin{bcode}
\begin{lstlisting}
q, k, v = self.qkv(self.ln1(x)).split(C, dim=2)
q, k, v = (t.view(B, T, heads, C // heads).transpose(1, 2) for t in (q, k, v))
att = (q @ k.transpose(-2, -1)) / math.sqrt(C // heads)
att = F.softmax(att.masked_fill(self.mask[:T, :T], float("-inf")), dim=-1)
x = x + self.proj((att @ v).transpose(1, 2).reshape(B, T, C))
x = x + self.mlp(self.ln2(x))
\end{lstlisting}
\end{bcode}
\begin{bres}
MiniGPT: 4 blocks, 4 heads, 128-dimensional embeddings, 128-character context --- \textbf{826,433 parameters}. Validation loss falls from 3.596 to \textbf{1.642} nats per character in 1,200 AdamW steps: perplexity \textbf{5.16} versus 65 for random guessing. Samples look like Shakespeare in form --- speaker names, line breaks, plausible words.
\end{bres}
\figw[0.6]{s3_gpt_loss.png}{MiniGPT pre-training: train and validation loss stay close, so the model is not overfitting yet.}
\begin{bapp}
The same architecture, scaled to billions of parameters and trillions of tokens, powers ChatGPT, Claude, Gemini and Llama, and code completion in tools such as GitHub Copilot.
\end{bapp}

%% =====================================================================
\hourhead{12}{TEACH}{Encoders, transfer learning and baselines}{SST-2 (6,920 train / 1,821 test reviews)}{Transformer encoders, fine-tuning, TF-IDF}

\topic{Encoders, decoders and transfer learning}
\begin{bdef}[Encoder versus decoder]
An \textbf{encoder} (BERT-style) reads the whole text in both directions and is best for classification and search. A \textbf{decoder} (GPT-style) reads left to right and is best for generation.
\end{bdef}
\begin{bdef}[Transfer learning]
Pretrain a model on huge amounts of unlabelled text, then fine-tune it on a small labelled task. The pretrained model already knows language; fine-tuning teaches the task.
\end{bdef}
\begin{bmath}[TF-IDF baseline]
$\mathrm{tfidf}(t,d)=\mathrm{tf}(t,d)\cdot\log\dfrac{N}{\mathrm{df}(t)}$ --- words frequent in this review but rare across reviews get high weight.
\end{bmath}
\begin{bex}[Classifying with an encoder]
Prepend a \texttt{[CLS]} token, run the encoder, and feed the final \texttt{[CLS]} vector into a linear layer that outputs \emph{positive} or \emph{negative}.
\end{bex}
\begin{bcode}
\begin{lstlisting}
tfidf = TfidfVectorizer(ngram_range=(1, 2), min_df=2, sublinear_tf=True)
lr = LogisticRegression(max_iter=2000, C=4).fit(tfidf.fit_transform(train.text), train.label)
# In class on a GPU, a pretrained encoder:
model = AutoModelForSequenceClassification.from_pretrained("distilbert-base-uncased", num_labels=2)
\end{lstlisting}
\end{bcode}
\begin{bres}[SST-2 test accuracy]
TF-IDF + logistic regression: \textbf{0.810}. \qquad Transformer encoder trained from scratch: 0.749.
\end{bres}
\begin{bwarn}[Always beat a simple baseline first]
A transformer trained from scratch on 7,000 sentences loses to TF-IDF. Transformers win through \textbf{pretraining}; a pretrained DistilBERT fine-tuned for two epochs typically reaches about 0.91 on SST-2.
\end{bwarn}
\begin{bapp}
Review and social-media sentiment monitoring, support-ticket routing, spam and intent detection. Hugging Face hosts thousands of pretrained encoders ready to fine-tune.
\end{bapp}

%% =====================================================================
\hourhead{13}{BUILD}{Instruction tuning with LoRA}{SST-2 reviews in an instruction template}{PEFT, LoRA, chat templates, memory budgeting}

\topic{Instruction tuning and LoRA}
\begin{bdef}[Instruction tuning]
Fine-tuning a pretrained model on (instruction, response) pairs so it follows a task format instead of just continuing text.
\end{bdef}
\begin{bdef}[LoRA --- low-rank adaptation]
Freeze the pretrained weights and learn a small low-rank correction for selected layers. Only the correction is trained and saved.
\end{bdef}
\begin{bmath}[LoRA update]
$h=Wx+\dfrac{\alpha}{r}\,BAx$, \;$W\in\mathbb{R}^{d\times k}$ frozen, $A\in\mathbb{R}^{r\times k}$, $B\in\mathbb{R}^{d\times r}$, $r\ll\min(d,k)$, $B$ initialised to 0.\\[3pt]
Trainable parameters per layer: $r(d+k)$ instead of $dk$.
\end{bmath}
Because $B$ starts at zero, training starts exactly from the pretrained model. The rank $r$ controls capacity: small $r$ is cheap and resists forgetting; larger $r$ can learn more.
\begin{bex}[Worked example]
For a $4096\times4096$ layer, full fine-tuning updates $16.8$\,M weights; LoRA with $r=8$ updates $8\cdot(4096+4096)=65{,}536$ --- about 0.4\%.
\end{bex}
\begin{bdef}[QLoRA]
LoRA on top of a base model stored in 4-bit precision, which lets multi-billion-parameter models be fine-tuned on a single GPU.
\end{bdef}

\topic{Chat template and answer-only loss}
\begin{bcode}[Chat template]
\begin{lstlisting}
### Review:
a stirring , funny and finally transporting re imagining of beauty and the beast
### Sentiment: positive
\end{lstlisting}
\end{bcode}
\begin{balgo}[LoRA instruction tuning]
\begin{algorithmic}[1]
\State freeze all base weights
\State wrap attention and MLP layers with LoRA matrices $A,B$
\For{each batch of templated examples}
  \State set the targets of prompt tokens to $-100$ (ignored by the loss)
  \State compute the loss on the answer tokens only
  \State update $A$ and $B$
\EndFor
\end{algorithmic}
\end{balgo}
\begin{bcode}
\begin{lstlisting}
class LoRALinear(nn.Module):
    def __init__(self, base: nn.Linear, r=8, alpha=16):
        super().__init__(); self.base, self.scale = base, alpha / r
        self.A = nn.Parameter(torch.randn(r, base.in_features) * 0.01)
        self.B = nn.Parameter(torch.zeros(base.out_features, r))
    def forward(self, x):
        return self.base(x) + (x @ self.A.T @ self.B.T) * self.scale
\end{lstlisting}
\end{bcode}
\begin{bres}
Trainable LoRA parameters: \textbf{65,536 of 891,969 (7.3\%)}.
\end{bres}
\begin{bmath}[GPU-memory rule of thumb with AdamW]
Full fine-tuning needs about 16 bytes per parameter (weights, gradients and two Adam moments). For a 0.5\,B model: about 8\,GB, versus about 1\,GB of fp16 weights plus small adapters with LoRA.
\end{bmath}
\begin{bwarn}[Our first attempt failed --- and why]
We first trained LoRA on whole text windows, so almost all the learning signal went into predicting review text. Result: 48\% accuracy, worse than chance, with every review labelled ``negative''. Computing the loss only on the answer tokens fixed it --- which is exactly why instruction-tuning libraries mask the prompt.
\end{bwarn}
\begin{bapp}
Domain assistants for legal, medical and customer-support work are built by adapting open models this way. Tools: Hugging Face PEFT, Unsloth, Axolotl.
\end{bapp}

%% =====================================================================
\hourhead{14}{BUILD}{Evaluating language models}{SST-2 test (held out) + Shakespeare validation}{LLM evaluation, benchmark hygiene, contamination}

\topic{How to score an LLM}
\begin{bdef}[Likelihood scoring]
For a multiple-choice task, compute how probable the model finds each candidate answer after the prompt, and pick the most probable.
\end{bdef}
\begin{bmath}[Score]
$\hat a=\arg\max_{a\in\{\text{positive},\text{negative}\}}\;\sum_{t\in a}\log P(t\mid\text{prompt},\,a_{<t})$
\end{bmath}
\begin{bdef}[Contamination]
Test items, or near-copies, that appear in the training data. Scores then measure memory rather than skill.
\end{bdef}
\begin{balgo}[Evaluation harness]
\begin{algorithmic}[1]
\State freeze the test set and the prompts
\State run every system with identical settings
\State include simple baselines
\State check train--test overlap (contamination)
\State check forgetting on the original task
\State repeat with several random seeds and report the spread
\end{algorithmic}
\end{balgo}
\begin{bres}[SST-2 test accuracy (400 held-out reviews for the language models)]
\begin{tabular}{lr}
\toprule
System & Accuracy\\
\midrule
Base MiniGPT, zero-shot & 0.520\\
MiniGPT + LoRA (ours) & 0.605\\
Transformer encoder from scratch & 0.749\\
TF-IDF + logistic regression & \textbf{0.810}\\
\bottomrule
\end{tabular}\\[3pt]
Forgetting: Shakespeare validation loss 1.642 for the base model, 2.721 with the LoRA adapter active. Contamination: 0.3\% of test reviews share a 30-character chunk with the training text.
\end{bres}
LoRA lifted a character-level model trained on plays from chance to clearly above chance on sentiment, by training 7\% of its weights --- a real effect. But a small model is still a small model, and the simple baseline wins. With the adapter active the model is worse at Shakespeare; switching the adapter off restores the original model exactly, because the base weights never changed.
\begin{bdef}[LLM-as-judge]
A stronger model grades answers against a rubric. It is cheap and scalable but has biases: \emph{position bias} (prefers the first answer), \emph{length bias} and \emph{self-preference}. Swap the answer order and average.
\end{bdef}
\begin{bapp}
EleutherAI's \texttt{lm-evaluation-harness}, Stanford HELM and LMArena (crowd-sourced pairwise ratings) are widely used to compare models.
\end{bapp}

%% =====================================================================
\hourhead{15}{BUILD}{Decoding, structured output and serving}{Team prompt suite}{Inference serving, structured outputs, latency and cost}

\topic{Decoding strategies}
\begin{bdef}[Decoding]
The rule for choosing the next token from the model's probabilities: greedy (always the most likely), random sampling, top-$k$ or top-$p$ sampling.
\end{bdef}
\begin{bmath}[Temperature, top-$k$ and top-$p$]
Temperature: $p_i\propto e^{z_i/T}$ --- low $T$ sharpens (safe, repetitive), high $T$ flattens (creative, error-prone).\\
Top-$k$: keep only the $k$ most likely tokens. \;Top-$p$ (nucleus): keep the smallest set whose total probability is at least $p$.
\end{bmath}
\begin{bex}[Logits $(2,1,0)$]
At $T=1$: probabilities $(0.67, 0.24, 0.09)$. At $T=0.5$: $(0.87, 0.12, 0.02)$ --- the top choice dominates.
\end{bex}
\begin{bres}
A $T=0.8$ sample after ``KING:'' --- ``'Tis not to tell the Lord the night way my hand''. At $T=1.5$ the text breaks into invented words.
\end{bres}

\topic{Structured output and throughput}
\begin{bdef}[Structured output]
Constraining the answer to a schema --- JSON with fixed fields and allowed values --- so that software can rely on it instead of parsing free text.
\end{bdef}
\begin{bcode}
\begin{lstlisting}
{"review": "one of the best movies of the year", "sentiment": "positive", "confidence": 0.51}
\end{lstlisting}
\end{bcode}
\begin{bwarn}[Confidence is not correctness]
In our prompt suite the LoRA model labels every example ``positive'', even ``a dull, lifeless and forgettable film''. A 60\%-accurate model is often confidently wrong; route low-confidence cases to a human or a stronger model.
\end{bwarn}
\begin{bmath}[Throughput and the KV-cache]
Throughput $=\dfrac{\text{batch size}\times\text{tokens generated}}{\text{time}}$. A KV-cache stores the keys and values of earlier tokens so each new step computes only the new token.
\end{bmath}
\begin{bres}[Generation speed on a laptop CPU]
\begin{tabular}{rrr}
\toprule
Batch & ms per step & tokens per second\\
\midrule
1 & 1.76 & 569\\
4 & 3.68 & 1,086\\
16 & 7.74 & 2,066\\
\bottomrule
\end{tabular}\\[3pt]
Batching 16 requests raised throughput about $3.6\times$ while per-step latency rose about $4.4\times$ --- far less than $16\times$.
\end{bres}
\begin{bapp}
Serving engines: vLLM (continuous batching and paged KV-cache), Ollama for laptops, Hugging Face TGI. Commercial LLM APIs offer JSON or structured-output modes.
\end{bapp}

%% =====================================================================
\section*{Formula sheet}
\begin{tabularx}{\linewidth}{@{}l X@{}}
\toprule
Language model & $P(x_{1:T})=\prod_tP(x_t\mid x_{<t})$; \; perplexity $=e^{\mathcal{L}}$\\
Attention & $\mathrm{softmax}(QK^\top/\sqrt{d_k}+M)V$\\
Block & $h\leftarrow h+\mathrm{Attn}(\mathrm{LN}(h))$; \; $h\leftarrow h+\mathrm{MLP}(\mathrm{LN}(h))$\\
TF-IDF & $\mathrm{tf}\cdot\log(N/\mathrm{df})$\\
LoRA & $h=Wx+\frac{\alpha}{r}BAx$; \; trainable $r(d+k)$\\
Temperature & $p_i\propto e^{z_i/T}$\\
Throughput & batch $\times$ tokens / time\\
\bottomrule
\end{tabularx}

\section*{Check yourself}
\begin{enumerate}
\item Validation loss is 1.642 nats per character. Perplexity? \emph{$e^{1.642}\approx5.2$.}
\item Why divide attention scores by $\sqrt{d_k}$? \emph{Large dot products saturate the softmax and kill gradients.}
\item Trainable parameters for LoRA $r=16$ on a $1024\times1024$ layer? \emph{$16\times2048=32{,}768$.}
\item Why mask prompt tokens during instruction tuning? \emph{Otherwise the loss is dominated by predicting the prompt, not the answer.}
\item What does a 0.3\% train--test overlap tell you? \emph{Contamination is negligible; scores are not inflated by memory.}
\end{enumerate}

\section*{Glossary}
\textbf{Token} unit of text $\cdot$ \textbf{Embedding} learned vector $\cdot$ \textbf{Attention} weighted mixing of tokens $\cdot$ \textbf{Perplexity} $e^{\text{loss}}$ $\cdot$ \textbf{LoRA} low-rank adapter $\cdot$ \textbf{Instruction tuning} task-format fine-tuning $\cdot$ \textbf{Contamination} test data in training $\cdot$ \textbf{KV-cache} stored keys and values.

\begin{submitbox}
Deliverable \textbf{P3 --- MiniLLM}: notebook, LoRA adapter weights, evaluation table (with baselines and the forgetting check) and a JSON-serving example in \texttt{P3-minillm/}. Tag \texttt{sunday-3}, create a Release, submit the link.
\end{submitbox}

\clearpage
\section*{Monday--Saturday Project Homework: Domain Instruction-Tuned LLM}
\addcontentsline{toc}{section}{Monday--Saturday Project Homework}
\begin{tcolorbox}[enhanced,colback=navy,colframe=gold,boxrule=0.8pt,arc=3pt,coltext=white,left=8pt,right=8pt,top=6pt,bottom=6pt]
{\color{gold}\bfseries Weekly project.} Organisations want small, private models that follow their instructions well. You will adapt a real open model with LoRA, evaluate it rigorously and serve it.\\[3pt]
{\color{gold}\bfseries Data:} databricks-dolly-15k (instructions) + a small domain corpus you choose (e.g. your university handbook)\\[3pt]
{\color{gold}\bfseries Rhythm:} about 2--3 hours a day. Commit every day; tag \texttt{week-3} and create a GitHub Release by \textbf{Saturday 23:59 local time}. Mentors review it before the next Sunday.
\end{tcolorbox}

\begin{tcolorbox}[enhanced,breakable,colback=white,colframe=sun,boxrule=0.8pt,arc=2pt,left=7pt,right=7pt,top=4pt,bottom=4pt,
  title={\bfseries Monday --- Tokenisers in practice\hfill\normalfont\small 2 h},coltitle=white,colbacktitle=sun,fonttitle=\small]
\textbf{Goal.} See how tokenisation shapes cost and quality.
\begin{enumerate}[itemsep=1pt,topsep=2pt]
  \item Train a BPE tokenizer (Hugging Face \texttt{tokenizers}, 8k vocabulary) on your domain corpus.
  \item Compare tokens per sentence against the GPT-2 and Qwen2.5 tokenizers on 200 sentences, including Indian-language text.
  \item Explain why token counts matter for cost and context length.
\end{enumerate}
\textbf{\color{dataC}Deliverable:} \texttt{notebooks/01\_tokenizers.ipynb}, comparison table\\
\textbf{\color{navy}Done when:} Tokenizer trained and saved; counts reported for all three tokenizers with examples of splits.\\
\textbf{\color{skillC}Stretch:} Measure how many tokens a Hindi or Telugu sentence needs compared with its English translation.
\end{tcolorbox}

\begin{tcolorbox}[enhanced,breakable,colback=white,colframe=sun,boxrule=0.8pt,arc=2pt,left=7pt,right=7pt,top=4pt,bottom=4pt,
  title={\bfseries Tuesday --- Scale up MiniGPT\hfill\normalfont\small 3 h},coltitle=white,colbacktitle=sun,fonttitle=\small]
\textbf{Goal.} Turn the Sunday model into a stronger small LM.
\begin{enumerate}[itemsep=1pt,topsep=2pt]
  \item Switch MiniGPT from characters to your BPE tokens; raise context to 256.
  \item Add learning-rate warm-up, cosine decay and gradient clipping; train on a GPU.
  \item Compare validation perplexity of character-level and BPE models fairly (convert to bits per character).
\end{enumerate}
\textbf{\color{dataC}Deliverable:} Training script, loss curves, comparison table\\
\textbf{\color{navy}Done when:} Bits-per-character used for a fair comparison; samples shown at two temperatures.\\
\textbf{\color{skillC}Stretch:} Add rotary position embeddings and measure the effect.
\end{tcolorbox}

\begin{tcolorbox}[enhanced,breakable,colback=white,colframe=sun,boxrule=0.8pt,arc=2pt,left=7pt,right=7pt,top=4pt,bottom=4pt,
  title={\bfseries Wednesday --- LoRA / QLoRA on a real open model\hfill\normalfont\small 3 h},coltitle=white,colbacktitle=sun,fonttitle=\small]
\textbf{Goal.} Adapt a 0.5B-parameter model on one GPU.
\begin{enumerate}[itemsep=1pt,topsep=2pt]
  \item Load \texttt{Qwen2.5-0.5B-Instruct} (or \texttt{SmolLM2-360M-Instruct}) with \texttt{transformers}.
  \item Fine-tune with \texttt{peft} LoRA (r = 8 and r = 32) on a 2,000-example Dolly subset, masking the prompt tokens.
  \item Track loss, GPU memory and time; save the adapters.
\end{enumerate}
\textbf{\color{dataC}Deliverable:} Adapters on the Hugging Face Hub (private is fine) and \texttt{reports/finetune.md}\\
\textbf{\color{navy}Done when:} Prompt masking verified; both ranks compared on validation loss and memory.\\
\textbf{\color{skillC}Stretch:} Run QLoRA (4-bit) and compare memory and quality.
\end{tcolorbox}

\begin{tcolorbox}[enhanced,breakable,colback=white,colframe=sun,boxrule=0.8pt,arc=2pt,left=7pt,right=7pt,top=4pt,bottom=4pt,
  title={\bfseries Thursday --- A rigorous evaluation harness\hfill\normalfont\small 3 h},coltitle=white,colbacktitle=sun,fonttitle=\small]
\textbf{Goal.} Know whether the fine-tune actually helped.
\begin{enumerate}[itemsep=1pt,topsep=2pt]
  \item Hold out 200 Dolly items; compare base vs LoRA with an LLM-as-judge, swapping answer order to remove position bias.
  \item Run \texttt{lm-evaluation-harness} on two tasks before and after fine-tuning to check for forgetting.
  \item Check n-gram overlap between your training and test items (contamination).
\end{enumerate}
\textbf{\color{dataC}Deliverable:} \texttt{reports/evaluation.md} with win-rate, benchmark and overlap tables\\
\textbf{\color{navy}Done when:} Order-swapped judging reported; forgetting measured; contamination check documented.\\
\textbf{\color{skillC}Stretch:} Have two humans rate 30 answers and compare with the judge.
\end{tcolorbox}

\begin{tcolorbox}[enhanced,breakable,colback=white,colframe=sun,boxrule=0.8pt,arc=2pt,left=7pt,right=7pt,top=4pt,bottom=4pt,
  title={\bfseries Friday --- Serving, structured output and throughput\hfill\normalfont\small 2--3 h},coltitle=white,colbacktitle=sun,fonttitle=\small]
\textbf{Goal.} Run the model like a product.
\begin{enumerate}[itemsep=1pt,topsep=2pt]
  \item Serve the merged model with vLLM or Ollama behind an OpenAI-compatible API.
  \item Force JSON output with a schema for one task (e.g. extracting fields from a notice).
  \item Benchmark latency and tokens/second at batch sizes 1, 4, 16 and 32.
\end{enumerate}
\textbf{\color{dataC}Deliverable:} Serving config, JSON-schema demo, benchmark table\\
\textbf{\color{navy}Done when:} Invalid JSON rate measured; throughput curve plotted; hardware stated.\\
\textbf{\color{skillC}Stretch:} Compare 4-bit and 16-bit serving speed and quality.
\end{tcolorbox}

\begin{tcolorbox}[enhanced,breakable,colback=white,colframe=sun,boxrule=0.8pt,arc=2pt,left=7pt,right=7pt,top=4pt,bottom=4pt,
  title={\bfseries Saturday --- Chat UI, report and model card\hfill\normalfont\small 3 h},coltitle=white,colbacktitle=sun,fonttitle=\small]
\textbf{Goal.} Package and present.
\begin{enumerate}[itemsep=1pt,topsep=2pt]
  \item Build a small chat UI (Gradio) that calls your served model.
  \item Write \texttt{MODEL\_CARD.md} including evaluation, known failure cases and safety notes.
  \item Record a 5-minute demo; push, tag \texttt{week-3}, create a Release.
\end{enumerate}
\textbf{\color{dataC}Deliverable:} UI, model card, demo link, Release \texttt{week-3}\\
\textbf{\color{navy}Done when:} Model card lists at least 5 tested failure cases with examples.\\
\textbf{\color{skillC}Stretch:} Add a refusal policy for out-of-domain questions and test it.
\end{tcolorbox}


\end{document}
